\documentclass[12pt]{article}
\usepackage{fancyhdr} % 用于自定义页眉页脚
\usepackage{listings} % 用于代码高亮
\usepackage{xcolor} % 用于定义颜色
\usepackage{fontspec} % 用于设置字体 (需要XeLaTeX或LuaLaTeX)
\usepackage{graphicx} % 用于插入图片
\usepackage{float} % 用于强制图片位置
\usepackage{geometry} % 用于设置页面边距

\geometry{a4paper, margin=1in} % 设置A4纸大小和1英寸边距

% 设置图片路径
\graphicspath{{./fig/}} % 假设图片存放在images文件夹中

% 定义VS Code风格的颜色
\definecolor{vscode-bg}{RGB}{30,30,30}
\definecolor{vscode-string}{RGB}{206,145,120}
\definecolor{vscode-keyword}{RGB}{86,156,214}
\definecolor{vscode-comment}{RGB}{106,153,85}
\definecolor{vscode-number}{RGB}{181,206,168}
\definecolor{vscode-function}{RGB}{220,220,170}

% 设置代码块字体（Fira Code或Consolas或Monaco）
\setmonofont{Monaco}[Scale=0.9]

% 配置listings样式
\lstdefinestyle{vscode}{
    backgroundcolor=\color{vscode-bg},
    commentstyle=\color{vscode-comment},
    keywordstyle=\color{vscode-keyword}\bfseries,
    numberstyle=\tiny\color{white},
    stringstyle=\color{vscode-string},
    basicstyle=\ttfamily\footnotesize\color{white},
    breakatwhitespace=false,
    breaklines=true,
    captionpos=b,
    keepspaces=true,
    numbers=left,
    numbersep=5pt,
    showspaces=false,
    showstringspaces=false,
    showtabs=false,
    tabsize=2,
    frame=single,
    frameround=tttt,
    rulecolor=\color{gray!30}
}

% 设置默认代码样式
\lstset{style=vscode}

\title{GalPHOS Project Report}
\author{\textbf{Group 16}: Hengxu Wu, Kairui Li, Yangyi Xiong}
\date{}

% 设置页面样式
\pagestyle{fancy}
\fancyhf{} % 清除所有页眉页脚
\fancyhead[L]{\leftmark} % 左上方显示section
\fancyhead[R]{\rightmark} % 右上方显示subsection
\fancyfoot[C]{\thepage} % 页脚中央显示页码

% 重新定义section和subsection的标记格式
\renewcommand{\sectionmark}[1]{\markboth{#1}{}}
\renewcommand{\subsectionmark}[1]{\markright{#1}}

\begin{document}
\maketitle

\section{Introduction}
\subsection{Project Overview}
% 简单介绍项目的内容和情况

GalPHOS (Galaxy Physics Online System) is a physical competition event management system based on the web platform, covering a variety of functions including event registration, exam management, and score statistics. The system is designed to provide an efficient and convenient management platform for physics competitions, supporting a variety of user roles, including administrators, students, coaches and graders.

Our system architecture and functions are based on the life cycle design of a competition, including five stages: registration, pre-match preparation, answer submission, post-match marking, result statistics and archives. The figure below illustrates the workflow of the system.
\begin{figure}[H]
    \centering
    \includegraphics[width=\textwidth]{Fig1.png}
    \caption{GalPHOS System Workflow}
    \label{fig:workflow}
\end{figure}

\subsection{Our Work}
% 介绍小组成员的工作内容和分工

Our group consists of three members: Hengxu Wu, Kairui Li, and Yangyi Xiong. The leader of the group is Hengxu Wu. Each member has contributed to different aspects of the project:
\begin{itemize}
    \item Hengxu Wu: Responsible for project design, development documents, and front-end development.
    \item Kairui Li: % TODO
    \item Yangyi Xiong: % TODO
\end{itemize}

\section{System Architecture}
\subsection{Architecture Overview}
% 描述系统的整体架构和设计思路，包括前端和后端的分离
\subsection{Frontend Design}
% 介绍前端的设计和实现（类型这一块）
\subsection{Backend Design}
% 介绍后端的设计和实现（类型这一块）

\section{Usage and Features}
\subsection{Administrator}
\subsection{Student}
\subsection{Coach}
\subsection{Grader}

\section{Summary}
% 对项目的总结和展望
\end{document}